% =====================================================================
%  chapters/minggu-11.tex
%  Minggu 11 -- JavaScript: Event Handling & Manipulasi DOM
% =====================================================================

\chapter{JavaScript: Event Handling \& Manipulasi DOM}
\label{chap:mjs-11}

% ---------------------------------------------------------------------
% 1. CAPAIAN PEMBELAJARAN
% ---------------------------------------------------------------------
\begin{capaian}
Setelah menyelesaikan bab ini, mahasiswa diharapkan mampu:
\begin{itemize}
  \item merespons interaksi pengguna (klik, input, \emph{submit}) menggunakan \emph{event handling};
  \item membedakan penggunaan \texttt{onclick} dan \texttt{addEventListener};
  \item memanipulasi konten DOM (\texttt{textContent}, \texttt{innerHTML}, \texttt{style}, \texttt{classList});
  \item membuat dan menghapus elemen HTML secara dinamis.
\end{itemize}
\end{capaian}

% ---------------------------------------------------------------------
% 2. TUJUAN PRAKTIK
% ---------------------------------------------------------------------
\begin{tujuanpraktik}
Pada sesi praktik ini, mahasiswa akan:
\begin{itemize}
  \item membuat halaman web interaktif dengan tombol dan form;
  \item membangun daftar tugas yang bisa ditambah/dihapus;
  \item menerapkan \emph{toggle menu mobile} sebagai mini proyek.
\end{itemize}
\end{tujuanpraktik}

\begin{catatan}
Pengingat cepat dari Minggu 10: \texttt{getElementById}, \texttt{querySelector}, \texttt{querySelectorAll}, \texttt{let}, \texttt{const}, \texttt{if}/\texttt{else}, fungsi.
\end{catatan}

% =====================================================================
\section{Event Handling: onclick \& addEventListener}
\label{sec:11-event-basic}
% =====================================================================

\textbf{Event} adalah sesuatu yang terjadi pada halaman web---misal: pengguna mengklik tombol, mengetik di \emph{input}, atau menggerakkan \emph{mouse}. Kita bisa \textbf{mendengarkan} event ini dan menjalankan JavaScript sebagai respons.

\begin{langkah}
  \item Buat file baru \texttt{minggu11.html} dengan kode berikut.
  \item Buka dengan Live Server, klik kedua tombol, dan amati perubahan.
\end{langkah}

\begin{lstlisting}[language=HTML, caption={Event handling: onclick dan addEventListener}, label={lst:11-event-basic}]
<!DOCTYPE html>
<html lang="id">
<head>
    <meta charset="UTF-8">
    <meta name="viewport" content="width=device-width, initial-scale=1.0">
    <title>Minggu 11 -- Event Handling</title>
    <style>
        body { font-family: Arial, sans-serif; margin: 20px; }
        button { padding: 10px 20px; font-size: 16px;
                 cursor: pointer; margin: 5px; }
        .aktif { background-color: #4CAF50; color: white; }
        .output { background-color: #f0f0f0; padding: 15px;
                  margin: 10px 0; border-radius: 5px;
                  min-height: 30px; }
    </style>
</head>
<body>
    <h1>Event Handling</h1>

    <!-- Cara 1: onclick langsung di HTML -->
    <h2>Cara 1: onclick di HTML</h2>
    <button onclick="sapa()">Klik Saya!</button>

    <!-- Cara 2: addEventListener di JavaScript -->
    <h2>Cara 2: addEventListener</h2>
    <button id="tombol-hitung">Hitung Klik</button>
    <span id="tampil-klik">0</span> kali diklik

    <div class="output" id="area-output">
        Output akan muncul di sini...
    </div>

    <script>
        // Cara 1: onclick
        function sapa() {
            document.getElementById("area-output").textContent =
                "Halo! Anda baru saja mengklik tombol sapa.";
        }

        // Cara 2: addEventListener (lebih disarankan)
        let jumlahKlik = 0;
        document.getElementById("tombol-hitung")
            .addEventListener("click", function() {
                jumlahKlik++;
                document.getElementById("tampil-klik")
                    .textContent = jumlahKlik;

                if (jumlahKlik >= 5) {
                    document.getElementById("area-output")
                        .textContent = "Wah, Anda sudah mengklik "
                            + jumlahKlik + " kali!";
                } else {
                    document.getElementById("area-output")
                        .textContent = "Klik ke-" + jumlahKlik
                            + ". Coba terus!";
                }
            });
    </script>
</body>
</html>
\end{lstlisting}

\begin{tangkapanlayar}
  {assets/images/minggu-11/event-basic.png}
  {Halaman dengan tombol ``Hitung Klik'' yang sudah diklik beberapa kali dan angka penghitung.}
  {Buka \texttt{minggu11.html} di peramban, klik tombol beberapa kali, dan bandingkan.}
\end{tangkapanlayar}

\subsection{Mengapa addEventListener Lebih Baik?}

\begin{table}[htbp]
\centering
\caption{Perbandingan onclick dan addEventListener}
\label{tab:11-perbandingan-event}
\begin{tabular}{lll}
\toprule
\textbf{Aspek} & \textbf{onclick} & \textbf{addEventListener} \\
\midrule
Jumlah \emph{handler} & Hanya 1 & Banyak sekaligus \\
Pemisahan kode & JS tercampur di HTML & JS terpisah di \texttt{<script>} \\
Fleksibilitas & Terbatas & Bisa \emph{add}/\emph{remove listener} \\
\bottomrule
\end{tabular}
\end{table}

\begin{tip}
Gunakan \texttt{addEventListener} untuk semua kebutuhan \emph{event handling}.
\end{tip}

\subsection{Event Populer}

\begin{table}[htbp]
\centering
\caption{Event yang sering digunakan}
\label{tab:11-event-populer}
\begin{tabular}{ll}
\toprule
\textbf{Event} & \textbf{Kapan Terjadi} \\
\midrule
\texttt{click} & Elemen diklik \\
\texttt{dblclick} & Elemen diklik dua kali \\
\texttt{mouseover} & \emph{Mouse} masuk ke elemen \\
\texttt{mouseout} & \emph{Mouse} keluar dari elemen \\
\texttt{keydown} & Tombol \emph{keyboard} ditekan \\
\texttt{input} & Nilai \emph{input} berubah \\
\texttt{submit} & Form dikirim \\
\texttt{load} & Halaman selesai dimuat \\
\bottomrule
\end{tabular}
\end{table}

% =====================================================================
\section{Event pada Form \& Input}
\label{sec:11-form-event}
% =====================================================================

\begin{langkah}
  \item Tambahkan kode berikut \textbf{sebelum} \texttt{</body>} pada \texttt{minggu11.html} (ganti atau tambahkan \emph{script} baru).
  \item Coba isi form dan \emph{submit}. Amati: jika field kosong $\rightarrow$ muncul pesan error; saat mengetik $\rightarrow$ muncul penghitung karakter \emph{real-time}; setelah \emph{submit} berhasil $\rightarrow$ form terkosongkan.
\end{langkah}

\begin{lstlisting}[language=HTML, caption={Form sederhana dengan validasi dan penghitung karakter}, label={lst:11-form}]
<h2>Form Sederhana</h2>
<form id="form-mahasiswa">
    <label for="input-nama">Nama:</label><br>
    <input type="text" id="input-nama"
           placeholder="Masukkan nama"><br><br>
    <label for="input-email">Email:</label><br>
    <input type="email" id="input-email"
           placeholder="contoh@email.com"><br><br>
    <button type="submit">Kirim</button>
</form>

<div class="output" id="hasil-form">
    Data form akan muncul di sini...
</div>

<script>
    // Event pada form submit
    document.getElementById("form-mahasiswa")
        .addEventListener("submit", function(event) {
            event.preventDefault();

            const nama = document.getElementById("input-nama").value;
            const email = document.getElementById("input-email").value;

            if (nama === "" || email === "") {
                document.getElementById("hasil-form").textContent =
                    "Semua field wajib diisi!";
                document.getElementById("hasil-form")
                    .style.backgroundColor = "#ffe0e0";
                return;
            }

            document.getElementById("hasil-form").textContent =
                "Data diterima -- Nama: " + nama
                    + ", Email: " + email;
            document.getElementById("hasil-form")
                .style.backgroundColor = "#e0ffe0";

            // Kosongkan form
            document.getElementById("input-nama").value = "";
            document.getElementById("input-email").value = "";
        });

    // Event pada input -- real-time feedback
    document.getElementById("input-nama")
        .addEventListener("input", function(event) {
            const panjang = event.target.value.length;
            if (panjang > 0) {
                document.getElementById("hasil-form").textContent =
                    "Anda mengetik " + panjang + " karakter...";
                document.getElementById("hasil-form")
                    .style.backgroundColor = "#f0f0f0";
            } else {
                document.getElementById("hasil-form").textContent =
                    "Data form akan muncul di sini...";
            }
        });
</script>
\end{lstlisting}

\begin{tangkapanlayar}
  {assets/images/minggu-11/event-form.png}
  {Form yang sudah diisi dengan data valid dan pesan sukses di area \emph{output}.}
  {Isi form dengan nama dan email yang valid, lalu klik ``Kirim''.}
\end{tangkapanlayar}

\begin{tangkapanlayar}
  {assets/images/minggu-11/event-form-validasi.png}
  {Form yang dikosongkan lalu disubmit---muncul pesan validasi error.}
  {Kosongkan form lalu klik ``Kirim'' tanpa mengisi apapun.}
\end{tangkapanlayar}

% =====================================================================
\section{Manipulasi DOM: Mengubah Konten}
\label{sec:11-manipulasi-konten}
% =====================================================================

\begin{langkah}
  \item Tambahkan kode berikut ke dalam \texttt{<body>} \texttt{minggu11.html}.
  \item Klik setiap tombol satu per satu dan amati perubahan.
\end{langkah}

\begin{lstlisting}[language=HTML, caption={Empat cara mengubah konten DOM}, label={lst:11-manipulasi-konten}]
<h2>Mengubah Konten DOM</h2>
<button id="btn-ubah-teks">Ubah Teks</button>
<button id="btn-ubah-html">Ubah HTML</button>
<button id="btn-ubah-style">Ubah Gaya</button>
<button id="btn-ubah-class">Toggle Class</button>

<div id="target-konten" class="output">
    Ini adalah elemen yang akan diubah.
</div>

<script>
    // 1. Mengubah teks (textContent)
    document.getElementById("btn-ubah-teks")
        .addEventListener("click", function() {
            document.getElementById("target-konten")
                .textContent = "Teks sudah berubah!";
        });

    // 2. Mengubah HTML (innerHTML)
    document.getElementById("btn-ubah-html")
        .addEventListener("click", function() {
            document.getElementById("target-konten")
                .innerHTML =
                "<strong>Teks tebal</strong> dan "
                + "<em>text miring</em> "
                + "-- semua dari innerHTML!";
        });

    // 3. Mengubah gaya langsung (style)
    document.getElementById("btn-ubah-style")
        .addEventListener("click", function() {
            const el = document.getElementById("target-konten");
            el.style.backgroundColor = "#4CAF50";
            el.style.color = "white";
            el.style.padding = "20px";
            el.style.borderRadius = "10px";
            el.style.fontSize = "18px";
        });

    // 4. Toggle class (cara paling rapi untuk styling)
    document.getElementById("btn-ubah-class")
        .addEventListener("click", function() {
            document.getElementById("target-konten")
                .classList.toggle("aktif");
        });
</script>
\end{lstlisting}

\begin{tangkapanlayar}
  {assets/images/minggu-11/manipulasi-konten.png}
  {Area \emph{output} setelah diklik---tampilan berubah sesuai tombol yang ditekan.}
  {Klik setiap tombol satu per satu dan bandingkan perubahan tampilan.}
\end{tangkapanlayar}

\begin{table}[htbp]
\centering
\caption{Ringkasan cara mengubah DOM}
\label{tab:11-cara-dom}
\begin{tabular}{lll}
\toprule
\textbf{Metode} & \textbf{Fungsi} & \textbf{Contoh} \\
\midrule
\texttt{textContent} & Mengubah teks mentah & \texttt{el.textContent = "Halo"} \\
\texttt{innerHTML} & Mengubah isi HTML & \texttt{el.innerHTML = "<b>Halo</b>"} \\
\texttt{style.prop} & Mengubah gaya CSS & \texttt{el.style.color = "red"} \\
\texttt{classList.toggle()} & Tambah/hapus class & \texttt{el.classList.toggle("aktif")} \\
\bottomrule
\end{tabular}
\end{table}

\begin{peringatan}
Hindari \texttt{innerHTML} untuk konten dari pengguna (bisa menyebabkan \textbf{XSS}). Untuk teks biasa, gunakan \texttt{textContent}.
\end{peringatan}

% =====================================================================
\section{Manipulasi DOM: Membuat \& Menghapus Elemen}
\label{sec:11-manipulasi-elemen}
% =====================================================================

\begin{langkah}
  \item Tambahkan kode berikut ke dalam \texttt{<body>} \texttt{minggu11.html}.
  \item Coba tambahkan beberapa item, hapus satu per satu, dan hapus semua.
\end{langkah}

\begin{lstlisting}[language=HTML, caption={Membuat dan menghapus elemen secara dinamis}, label={lst:11-hapus-elemen}]
<h2>Membuat & Menghapus Elemen</h2>
<input type="text" id="input-item"
       placeholder="Ketik item baru...">
<button id="btn-tambah">Tambah</button>
<button id="btn-hapus-terakhir">Hapus Terakhir</button>
<button id="btn-hapus-semua">Hapus Semua</button>

<ul id="daftar-item" style="list-style: none; padding: 0;">
</ul>

<script>
    function buatItem(teksItem) {
        const li = document.createElement("li");
        li.textContent = teksItem;
        li.style.padding = "8px 12px";
        li.style.margin = "5px 0";
        li.style.backgroundColor = "#e8f4f8";
        li.style.borderRadius = "4px";
        li.style.display = "flex";
        li.style.justifyContent = "space-between";

        const tombolHapus = document.createElement("button");
        tombolHapus.textContent = "x";
        tombolHapus.style.backgroundColor = "#ff6b6b";
        tombolHapus.style.color = "white";
        tombolHapus.style.border = "none";
        tombolHapus.style.padding = "2px 8px";
        tombolHapus.style.cursor = "pointer";
        tombolHapus.style.borderRadius = "3px";

        tombolHapus.addEventListener("click", function() {
            li.remove();
        });

        li.appendChild(tombolHapus);
        return li;
    }

    // Tambah item
    document.getElementById("btn-tambah")
        .addEventListener("click", function() {
            const input = document.getElementById("input-item");
            const teks = input.value.trim();
            if (teks === "") {
                alert("Masukkan teks terlebih dahulu!");
                return;
            }
            document.getElementById("daftar-item")
                .appendChild(buatItem(teks));
            input.value = "";
            input.focus();
        });

    // Tambah item dengan Enter
    document.getElementById("input-item")
        .addEventListener("keydown", function(event) {
            if (event.key === "Enter") {
                document.getElementById("btn-tambah").click();
            }
        });

    // Hapus item terakhir
    document.getElementById("btn-hapus-terakhir")
        .addEventListener("click", function() {
            const daftar = document.getElementById("daftar-item");
            if (daftar.lastElementChild) {
                daftar.lastElementChild.remove();
            }
        });

    // Hapus semua item
    document.getElementById("btn-hapus-semua")
        .addEventListener("click", function() {
            document.getElementById("daftar-item")
                .innerHTML = "";
        });
</script>
\end{lstlisting}

\begin{tangkapanlayar}
  {assets/images/minggu-11/manipulasi-elemen.png}
  {Daftar item yang sudah ditambahkan beberapa item, dengan tombol hapus per item.}
  {Buka \texttt{minggu11.html} di peramban, tambahkan item baru, lalu hapus satu per satu.}
\end{tangkapanlayar}

\subsection{Cara Membuat \& Menambah Elemen}

\begin{lstlisting}[language=JavaScript, caption={Pola dasar pembuatan elemen DOM}, label={lst:11-pola-create}]
// 1. Buat elemen baru
const elemenBaru = document.createElement("div");

// 2. Isi konten
elemenBaru.textContent = "Ini elemen baru";
elemenBaru.className = "kartu";

// 3. Masukkan ke DOM (sebagai anak dari elemen lain)
document.getElementById("induk").appendChild(elemenBaru);

// Atau sisipkan di posisi tertentu
document.getElementById("induk")
    .insertBefore(elemenBaru, referensiElemen);

// Hapus elemen
elemenBaru.remove();
\end{lstlisting}

% =====================================================================
\section{Mini Proyek: Toggle Menu Mobile}
\label{sec:11-mini-proyek}
% =====================================================================

Membuat tombol \emph{hamburger} sederhana yang membuka/menutup menu navigasi---seperti yang sering Anda lihat di situs web pada tampilan \emph{mobile}.

\begin{langkah}
  \item Buat file baru \texttt{mini-proyek-menu.html}.
  \item Buka dengan Live Server, lalu \textbf{resize browser ke bawah 600px} untuk melihat tampilan \emph{mobile}.
  \item Klik tombol $\equiv$ dan amati menu terbuka. Klik $\times$ untuk menutup.
\end{langkah}

\begin{lstlisting}[language=HTML, caption={Mini proyek: toggle menu mobile}, label={lst:11-toggle-menu}]
<!DOCTYPE html>
<html lang="id">
<head>
    <meta charset="UTF-8">
    <meta name="viewport"
          content="width=device-width, initial-scale=1.0">
    <title>Mini Proyek -- Toggle Menu</title>
    <style>
        * { margin: 0; padding: 0; box-sizing: border-box; }
        body { font-family: Arial, sans-serif; }

        .navbar {
            background-color: #333;
            padding: 15px 20px;
            display: flex;
            justify-content: space-between;
            align-items: center;
        }
        .navbar .logo {
            color: white; font-size: 20px; font-weight: bold;
        }
        .hamburger {
            display: none; background: none; border: none;
            color: white; font-size: 24px; cursor: pointer;
            padding: 5px 10px;
        }
        .nav-menu {
            list-style: none;
            display: flex; gap: 15px;
        }
        .nav-menu a {
            color: white; text-decoration: none;
            padding: 5px 10px; border-radius: 3px;
        }
        .nav-menu a:hover { background-color: #555; }

        @media (max-width: 600px) {
            .hamburger { display: block; }
            .nav-menu {
                display: none;
                flex-direction: vertical;
                position: absolute; top: 60px;
                left: 0; right: 0;
                background-color: #444;
                padding: 10px 20px;
            }
            .nav-menu.tampil { display: flex; }
        }

        .konten { padding: 20px; }
    </style>
</head>
<body>
    <nav class="navbar">
        <div class="logo">MyWeb</div>
        <button class="hamburger" id="hamburger">
            &#9776;
        </button>
        <ul class="nav-menu" id="nav-menu">
            <li><a href="#">Beranda</a></li>
            <li><a href="#">Tentang</a></li>
            <li><a href="#">Portofolio</a></li>
            <li><a href="#">Kontak</a></li>
        </ul>
    </nav>

    <div class="konten">
        <h2>Selamat Datang!</h2>
        <p>Coba ubah ukuran browser ke bawah 600px
           (atau buka di hp) untuk melihat menu
           hamburger muncul.</p>
    </div>

    <script>
        const tombolHamburger =
            document.getElementById("hamburger");
        const navMenu = document.getElementById("nav-menu");

        tombolHamburger.addEventListener("click", function() {
            navMenu.classList.toggle("tampil");
            if (navMenu.classList.contains("tampil")) {
                tombolHamburger.textContent = "\u2715";
            } else {
                tombolHamburger.textContent = "\u2630";
            }
        });

        document.addEventListener("click", function(event) {
            if (!navMenu.contains(event.target)
                && !tombolHamburger.contains(event.target)) {
                navMenu.classList.remove("tampil");
                tombolHamburger.textContent = "\u2630";
            }
        });
    </script>
</body>
</html>
\end{lstlisting}

\begin{tangkapanlayar}
  {assets/images/minggu-11/miniproject-desktop.png}
  {Tampilan desktop---menu horizontal, tombol \emph{hamburger} tersembunyi.}
  {Buka \texttt{mini-proyek-menu.html} di peramban dengan lebar penuh.}
\end{tangkapanlayar}

\begin{tangkapanlayar}
  {assets/images/minggu-11/miniproject-mobile.png}
  {Tampilan \emph{mobile} (browser $<$ 600\,px)---tombol \emph{hamburger} terlihat, menu terbuka ke bawah.}
  {Resize browser ke bawah 600\,px, lalu klik tombol $\equiv$.}
\end{tangkapanlayar}

% =====================================================================
\section{Latihan}
\label{sec:11-latihan}
% =====================================================================

\begin{latihan}[Kalkulator Sederhana]
Buat file \texttt{latihan11-1.html} yang berisi:
\begin{enumerate}
  \item Dua \emph{input} angka (\texttt{<input type="number">}).
  \item Empat tombol: Tambah, Kurang, Kalikan, Bagi.
  \item Area \emph{output} untuk menampilkan hasil.
  \item JavaScript yang menghitung sesuai tombol yang diklik.
\end{enumerate}

\textbf{Output yang diharapkan:}
\begin{itemize}
  \item Jika \emph{input} 10 dan 5, klik ``Kalikan'' $\rightarrow$ \emph{output} menampilkan \texttt{Hasil: 50}.
  \item Jika \emph{input} 10 dan 0, klik ``Bagi'' $\rightarrow$ \emph{output} menampilkan \texttt{Error: Tidak bisa dibagi nol!}.
\end{itemize}

\begin{tangkapanlayar}
  {assets/images/minggu-11/latihan1-kalkulator.png}
  {Kalkulator berfungsi dengan hasil perhitungan.}
  {Buka \texttt{latihan11-1.html} di peramban, masukkan dua angka, dan klik salah satu tombol.}
\end{tangkapanlayar}
\end{latihan}

\begin{latihan}[Daftar Tugas dengan Checkbox]
Buat file \texttt{latihan11-2.html} yang berisi:
\begin{enumerate}
  \item \emph{Input} untuk menambah tugas baru + tombol ``Tambah''.
  \item Setiap tugas memiliki \emph{checkbox} dan tombol ``Hapus''.
  \item Saat \emph{checkbox} dicentang, teks tugas dicoret (\emph{strikethrough}).
  \item Total tugas dan tugas selesai ditampilkan di bawah daftar.
\end{enumerate}

\textbf{Output yang diharapkan:}
\begin{itemize}
  \item Menambah 3 tugas $\rightarrow$ \texttt{Total: 3, Selesai: 0}.
  \item Mencentang 1 tugas $\rightarrow$ \texttt{Total: 3, Selesai: 1}, teks tugas tercoret.
\end{itemize}

\begin{tangkapanlayar}
  {assets/images/minggu-11/latihan2-tugas.png}
  {Daftar tugas dengan satu tugas tercentang dan \emph{counter} yang benar.}
  {Buka \texttt{latihan11-2.html} di peramban, tambahkan 3 tugas, centang salah satu.}
\end{tangkapanlayar}
\end{latihan}

\begin{latihan}[Light/Dark Mode Toggle --- Opsional]
Buat file \texttt{latihan11-3.html} yang berisi:
\begin{enumerate}
  \item Tombol ``Dark Mode'' / ``Light Mode''.
  \item Saat diklik, halaman berganti antara \emph{background} terang dan gelap.
  \item Preferensi diingat selama sesi (boleh pakai \texttt{localStorage} jika sudah tahu, atau cukup \emph{toggle} CSS \emph{class}).
\end{enumerate}

\begin{tangkapanlayar}
  {assets/images/minggu-11/latihan3-darkmode.png}
  {Halaman dalam mode gelap dan mode terang.}
  {Buka \texttt{latihan11-3.html} di peramban dan klik tombol \emph{toggle}.}
\end{tangkapanlayar}
\end{latihan}

% =====================================================================
\section{Troubleshooting}
\label{sec:11-troubleshooting}
% =====================================================================

\begin{peringatan}
\textbf{Event tidak merespons} --- pastikan elemen sudah ada di DOM saat script dijalankan (pindahkan \texttt{<script>} ke akhir \texttt{<body>} atau pakai \texttt{DOMContentLoaded}). Periksa \emph{selector} ID/\emph{class} di DevTools dan ejaan nama \emph{event} (mis.\ \texttt{"click"}, bukan \texttt{"clicl"}).
\end{peringatan}

\begin{catatan}
\textbf{Form refresh saat dikirim} --- Anda lupa menambahkan \texttt{event.preventDefault()}. Bungkus \emph{handler} \texttt{submit} dengan:
\end{catatan}

\begin{lstlisting}[language=JavaScript, label={lst:11-prevent-default}]
form.addEventListener("submit", function(event) {
    event.preventDefault();
    // proses data...
});
\end{lstlisting}

\begin{catatan}
\textbf{Elemen baru tidak muncul di DOM} --- pastikan Anda memanggil \texttt{appendChild} atau \texttt{prepend} setelah \texttt{createElement}. Tanpa panggilan tersebut, elemen tidak akan muncul.
\end{catatan}

\begin{tip}
\textbf{classList.toggle tidak berfungsi} --- periksa ejaan nama \emph{class} (\texttt{classList} bersifat \emph{case-sensitive}). Juga pastikan pengecekan \texttt{event.target} pada toggle menu benar menggunakan \texttt{contains()}.
\end{tip}
